问题二更像是在搭建一个“珍稀物种响应模块”。题面附件中的事实需要逐一映射到参数：2022 年公报对应种群初值和监测基准，附件 1 中的航道阻隔、鱼道连通与栖息地破碎化对应 \(\delta_B\)，附件 2 中的放流、补充和食源恢复对应 \(R_S\)，\(e_D,m_D,H,u\) 则分别反映繁殖效率、自然死亡、捕捞/扰动和污染压力。只有将这些事实与参数对应起来，才能分析禁渔收益如何在食物网和生境连通性上被放大或削弱，并解释同一背景下江豚和长江鲟不同的恢复曲线\cite{report2022}。

题面中的事实映射还意味着，问题二必须把“自然恢复”和“工程补偿”分开解释。江豚主要体现食源改善后的数量响应；对于长江鲟，除食源外，还要看洄游通道是否连续、放流是否提供额外补充，以及污染是否抵消恢复收益。因此，本问的重点不是比较哪个物种的绝对值更大，而是考察政策收益在不同约束条件下保留和损失的程度。明确这层关系，问题二才能由“珍稀物种是否增加”转向“珍稀物种响应机制如何分化”。

各参数都对应具体的事实链，而不是把变量当成脱离附件事实的纯数学符号。问题二的证据关系由禁渔和食源恢复出发，经通道、放流与污染修正，最终落到珍稀物种的数量变化。
